% Original last page: the frontiers collection.
\sectionlink*{The Frontiers Collection}
シリーズ編集者：A. C. Elitzur、M. P. Silverman、J. Tuszynski、R. Vaas、H. D. Zeh
\begin{description}
\item[相対性と時空の本性] V. Petkov著。
\item[量子力学はどこへ向かうのか] A. C. Elitzur、S. Dolev、N. Kolenda編。
\item[生命：脂質の問題として] 新たに現れる脂質科学、リピドミクス。O. G. Mouritsen著。
\item[量子と古典の類似] D. Dragoman、M. Dragoman著。
\item[知識と世界] サイエンス・ウォーズを超える課題。M. Carrier、J. Roggenhofer、G. Küppers、P. Blanchard編。
\item[量子と古典の対応] A. O. Bolivar著。
\item[心、物質、量子力学] H. Stapp著。
\item[量子力学と重力] M. Sachs著。
\item[自然と社会における極端な事象] S. Albeverio、V. Jentsch、H. Kantz編。
\item[生命の熱力学的機構] M. Kurzynski著。
\item[新たに現れる意識の物理学] J. A. Tuszynski編。
\item[弱い結び付き] タンパク質から社会ネットワークまで、複雑系の安定化要因。P. Csermely著。
\item[心、物質、内在秩序] P. T. I. Pylkkänen著。
\item[岐路に立つ量子力学] 歴史、哲学、物理学からの新たな視点。J. Evans、A. S. Thorndike編。
\item[粒子の形而上学] 亜原子的実在の批判的考察。B. Falkenburg著。
\item[時間の向きの物理的基礎] H. D. Zeh著。
\item[非対称性：情報の基礎] S. J. Muller著。
\item[心ある宇宙] 量子力学と参加する観測者。H. Stapp著。
\item[デコヒーレンスと量子から古典への移行] M. Schlosshauer著。
\item[量子重ね合わせ] コヒーレンス、もつれ、干渉の直観に反する帰結。Mark P. Silverman著。
\item[非線形の宇宙] カオス、創発、生命。A. Scott著。
\item[対称性が支配する] 科学と自然はいかに対称性に基づいているか。J. Rosen著。
\end{description}
